\documentclass[11pt]{article}
\usepackage[margin=1in]{geometry}
\usepackage{amsmath,amssymb,amsthm,mathtools}
\usepackage{graphicx}
\usepackage{microtype}
\usepackage[hidelinks]{hyperref}
\usepackage{enumitem}

\newtheorem{theorem}{Theorem}
\newtheorem{lemma}{Lemma}
\newcommand{\E}{\mathsf E}
\newcommand{\1}{\mathbf 1}
\newcommand{\Hatb}{H^{p}_{\mathrm{atb},q}}

\title{Both \(p<1\) Atomic-Block Questions Fail\\
on an Atomic Filtration}
\author{Solution packet for arXiv:1407.5451, Problem 3.1}
\date{22 July 2026}

\begin{document}
\maketitle

\begin{center}
\fbox{\parbox{0.91\textwidth}{\textbf{Status.}
Counterexample likely valid, subject to human review.  For every
\(0<p<1\) and \(1<q<\infty\), both assertions in the source problem fail
on the same countably atomic probability space with one fixed filtration.}}
\end{center}

\section{Source problem}

Conde-Alonso and Parcet define \((p,q)\)-atomic blocks for \(0<p<1\),
then ask in Problem 3.1 whether
\[
 H^p_{\mathrm{atb},q}(\Omega)=H_p(\Omega)
 \qquad\text{and}\qquad
 H_p(\Omega)^*=\Lambda_{p,q}(\Omega).
\]
Here \(H_p\) is defined by the martingale square function.  The proposed
Lipschitz norm is
\begin{equation}\label{eq:lambda}
 \|h\|_{\Lambda_{p,q}}
 =\sup_{k\geq1,\,A\in\Sigma_k}
 \frac{1}{\mu(A)^{1/p-1}}
 \left[
   \left(\frac1{\mu(A)}\int_A |h-\E_kh|^q\,d\mu\right)^{1/q}
   +\|dh_k\|_\infty
 \right].
\end{equation}
The question and the definitions immediately preceding it occur on page 17
of the source paper \cite{source}.

\begin{figure}[ht]
\centering
\includegraphics[width=0.94\textwidth]{figures/open_problem_crop.png}
\caption{The definition and Problem 3.1 on page 17 of arXiv:1407.5451.}
\end{figure}

For a \((p,q)\)-subatom supported on \(A_j\in\Sigma_{k_j}\), the source
requires
\begin{equation}\label{eq:subatom}
 \|a_j\|_q
 \leq \mu(A_j)^{1-1/p-1/q'}\frac1{k_j-k+1}
 =\mu(A_j)^{1/q-1/p}\frac1{k_j-k+1}.
\end{equation}
The size of one block \(b=\sum_j\lambda_ja_j\) is the infimum of
\(\sum_j|\lambda_j|\), while the outer decomposition of a function into
blocks is aggregated in \(\ell_p\).

\section{Proof intuition}

The definition has an exponent mismatch in how it groups pieces.  Inside one
block, arbitrarily many subatoms are charged in \(\ell_1\).  For \(p<1\), a
martingale square function can see the same pieces in \(\ell_p\), which is
much larger.

The concrete gadget has two signs.  At level two it equals \(+1\) or \(-1\).
At level three it usually cancels back to zero, but with probability
\(\varepsilon\) it jumps to size \(1/\varepsilon\) in order to preserve the
martingale condition.  Its terminal \(L_p\) norm tends to zero like
\(\varepsilon^{1/p-1}\), and the proposed atomic-block norm has the same
small upper bound.  Its square function, however, remembers the level-two
jump of size one on the whole block.  Countably many disjoint copies give a
function in \(L_p\cap H^p_{\mathrm{atb},q}\) but not in \(H_p\).

For the dual assertion, use a decaying sum of the level signs.  Its decay is
chosen so that testing every martingale difference costs at most its local
\(L_p\) square-function mass.  Formula \eqref{eq:lambda} nevertheless assigns
it infinite norm: at late levels its nonzero global jump is multiplied by a
negative power of the measure of an already resolved, extremely small atom.

\section{The counterexample}

\begin{theorem}\label{thm:main}
Fix \(0<p<1\) and \(1<q<\infty\).  There is a countably atomic probability
space with a fixed filtration for which
\begin{enumerate}[label=\textup{(\roman*)}]
\item \(H^p_{\mathrm{atb},q}(\Omega)\not=H_p(\Omega)\), even as sets;
\item \(H_p(\Omega)^*\not=\Lambda_{p,q}(\Omega)\), with
      \(\Lambda_{p,q}\) defined by \eqref{eq:lambda}.
\end{enumerate}
\end{theorem}

\begin{proof}
Put
\[
 w_n=2^{-n},\qquad
 \varepsilon_n=2^{-n/(1-p)} \quad(n\geq1).
\]
Thus \(\sum_nw_n=1\), \(0<\varepsilon_n\leq1/2\), and
\begin{equation}\label{eq:epssum}
 \varepsilon_n^{1-p}=2^{-n}.
\end{equation}
Partition \(\Omega\) into atoms
\(A_n^+,C_n^+,A_n^-,C_n^-\), with
\[
 \mu(A_n^\pm)=\frac{w_n\varepsilon_n}{2},
 \qquad
 \mu(C_n^\pm)=\frac{w_n(1-\varepsilon_n)}{2}.
\]
Let \(B_n^\pm=A_n^\pm\cup C_n^\pm\), put
\(B_n=B_n^+\cup B_n^-\), and let
\(R_n=\bigcup_{j\geq n}B_j\).  Define an increasing filtration by declaring
that the atoms of \(\Sigma_{2n-1}\) are
\[
 A_j^\pm,C_j^\pm\ (j<n),\qquad R_n,
\]
and that the atoms of \(\Sigma_{2n}\) are
\[
 A_j^\pm,C_j^\pm\ (j<n),\qquad B_n^+,B_n^-,R_{n+1}.
\]
Thus \(\Sigma_1\) is trivial.  From level \(2n-1\) to level \(2n\) the
filtration reveals the sign in the \(n\)-th block, and at level \(2n+1\) it
reveals whether the rare event occurred.

Define
\begin{equation}\label{eq:gn}
 g_n=\varepsilon_n^{-1}(\1_{A_n^+}-\1_{A_n^-}).
\end{equation}
The two signs have equal mass, so \(\E_{2n-1}g_n=0\), while
\[
 \E_{2n}g_n=\1_{B_n^+}-\1_{B_n^-}=:r_n,
 \qquad \E_{2n+1}g_n=g_n.
\]
Consequently the only nonzero differences are
\(dg_{n,2n}=r_n\) and \(dg_{n,2n+1}=g_n-r_n\).  The square
function \(S(g_n)\) is \(\sqrt2\) on \(C_n^+\cup C_n^-\), and it is
\[
 \left(1+\left(\frac{1-\varepsilon_n}{\varepsilon_n}\right)^2\right)^{1/2}
\]
on \(A_n^+\cup A_n^-\).  Therefore
\begin{align}
 \|g_n\|_p^p
   &=w_n\varepsilon_n^{1-p},\label{eq:lpgn}\\
 \|g_n\|_{H_p}^p
   &=w_n\left[(1-\varepsilon_n)2^{p/2}
      +\varepsilon_n
       \left(1+\left(\frac{1-\varepsilon_n}{\varepsilon_n}\right)^2
       \right)^{p/2}\right]
   \geq 2^{p/2-1}w_n.\label{eq:hpgn}
\end{align}

We next verify that \(g_n\) is a very cheap \((p,q)\)-atomic block.  Set
\(m_n=w_n\varepsilon_n/2\) and
\[
 a_n^+=\frac{m_n^{-1/p}}3\1_{A_n^+},
 \qquad
 a_n^-=-\frac{m_n^{-1/p}}3\1_{A_n^-},
 \qquad
 \lambda_n=\frac{3m_n^{1/p}}{\varepsilon_n}.
\]
Both subatoms are supported at level \(k_j=2n+1\), while the block has
cancellation at level \(k=2n-1\).  Moreover
\[
 \|a_n^\pm\|_q
 =\frac13m_n^{1/q-1/p}
 =m_n^{1/q-1/p}\frac1{(2n+1)-(2n-1)+1},
\]
so \eqref{eq:subatom} holds with equality, for every \(q>1\).  Since
\(g_n=\lambda_na_n^++\lambda_na_n^-\),
\begin{equation}\label{eq:blockgn}
 |g_n|_{\mathrm{atb},q}^p
 \leq2\lambda_n
 =6\,2^{-1/p}w_n^{1/p}\varepsilon_n^{1/p-1}.
\end{equation}

Let \(t_n=w_n^{-1/p}\).  The finite martingales
\begin{align*}
 F_{2N-1}&=\sum_{n<N}t_ng_n,\\
 F_{2N}&=\sum_{n<N}t_ng_n+t_Nr_N,\\
 F_{2N+1}&=\sum_{n\leq N}t_ng_n
\end{align*}
are integrable and compatible with the filtration.  They converge pointwise
and in \(L_p\) to
\[
 f=\sum_{n\geq1}t_ng_n.
\]
By \eqref{eq:lpgn} and disjointness,
\[
 \|f\|_p^p
 =\sum_n t_n^pw_n\varepsilon_n^{1-p}
 =\sum_n2^{-n}<\infty.
\]
Using the blocks \(t_ng_n\), equations \eqref{eq:blockgn} and
\eqref{eq:epssum} give
\[
 \|f\|_{H^p_{\mathrm{atb},q}}^p
 \leq (6\,2^{-1/p})^p
       \sum_n\varepsilon_n^{1-p}<\infty.
\]
Thus the series is an admissible atomic-block decomposition of the
\(L_p\)-function \(f\).  On the other hand, square functions are supported
on the disjoint sets \(B_n=B_n^+\cup B_n^-\).  Hence \eqref{eq:hpgn}
implies
\[
 \|f\|_{H_p}^p
 =\sum_n t_n^p\|g_n\|_{H_p}^p
 \geq 2^{p/2-1}\sum_n t_n^pw_n
 =2^{p/2-1}\sum_n1=\infty.
\]
This proves (i).

For (ii), put \(\alpha=1/p-1>0\),
\(\beta_n=(w_n/2)^\alpha\), and take the bounded function
\[
 \phi=\sum_{n\geq1}\beta_nr_n.
\]
Its only nonzero martingale differences are
\(d\phi_{2n}=\beta_nr_n\).  On finite martingales define
\[
 L_\phi(h)=\sum_{n\geq1}\beta_n\int (dh_{2n})r_n\,d\mu;
\]
this agrees with \(\int h\phi\,d\mu\) whenever the latter is directly
defined.  Write \(c_{n,\pm}\) for the two constant values of \(dh_{2n}\)
on \(B_n^\pm\), and put \(m_n'=w_n/2\).  Since
\(\beta_n=(m_n')^{1/p-1}\),
\begin{align*}
 \left|\beta_n\int (dh_{2n})r_n\,d\mu\right|
 &\leq \beta_nm_n'(|c_{n,+}|+|c_{n,-}|)\\
 &\leq
 \left(\int_{B_n}|dh_{2n}|^p\,d\mu\right)^{1/p}=:x_n.
\end{align*}
For \(p<1\), \(\sum_nx_n\leq(\sum_nx_n^p)^{1/p}\).  Therefore
\[
 |L_\phi(h)|\leq\left(\sum_n\int_{B_n}|dh_{2n}|^p\,d\mu\right)^{1/p}
 \leq\|h\|_{H_p}.
\]
Thus \(L_\phi\) extends continuously and \(\phi\) represents an element of
\(H_p^*\).

But \(\phi\notin\Lambda_{p,q}\).  In \eqref{eq:lambda}, take \(k=2n\)
and the already resolved atom \(A=A_{n-1}^+\), for \(n\geq2\).  Since
\(\|d\phi_{2n}\|_\infty=\beta_n\), the corresponding expression is at
least
\[
 \frac{\beta_n}{(w_{n-1}\varepsilon_{n-1}/2)^\alpha}
 =\left(\frac{1}{2\varepsilon_{n-1}}\right)^\alpha
 \longrightarrow\infty.
\]
This proves (ii).
\end{proof}

\section{Verification notes}

The proof uses only the definitions printed in the source.  In particular:
\begin{itemize}
\item the filtration is fixed once and for all; every finite stage of the
      separating martingale is integrable, and the finite stages converge in
      \(L_p\) to the displayed terminal function;
\item every subatom satisfies the source normalization with equality;
\item the separating function \(f\) belongs to \(L_p\), so the counterexample
      respects the set-level requirement in the definition of
      \(H^p_{\mathrm{atb},q}\);
\item the same construction works for every finite \(q>1\), because the
      \(q\)-dependence cancels in the indicator normalization;
\item the dual counterexample uses exactly the global jump term as it appears
      inside the brackets in the source's displayed \(\Lambda_{p,q}\) norm.
\end{itemize}

The accompanying script \texttt{code/verify\_counterexample.py} evaluates the
finite-block formulas and partial sums for several choices of \(p,q\).  It is a
sanity check, not a substitute for the proof.

\section{Novelty and limitations}

A bounded search was performed in the run indexes and on arXiv using
\texttt{1407.5451}, the exact notation \(H^p_{\mathrm{atb},q}\), the paper
title, and the phrases ``\((p,q)\)-atomic blocks'' and ``atomic block
\(H_p\) martingale''.  No later paper settling Problem 3.1 was found.  The
novelty claim is therefore bounded rather than exhaustive.

The theorem answers Problem 3.1 negatively for arbitrary filtrations.  It does
not rule out a repaired theorem under additional regularity, a restriction on
how subatoms may be grouped, or a modified Lipschitz norm in which the jump
term is not multiplied by the measure of an unrelated set.

\paragraph{Human-review recommendation.}
Verify the source transcription of the subatom exponent and the placement of
\(\|df_k\|_\infty\) in \eqref{eq:lambda}; then check the two conditional
expectations in \eqref{eq:gn}.  All remaining estimates are exact scalar
calculations.

\begin{thebibliography}{1}
\bibitem{source}
J.~M. Conde-Alonso and J. Parcet,
\emph{Atomic blocks for noncommutative martingales},
arXiv:1407.5451 (2014), Problem 3.1, p.~17.
\end{thebibliography}

\end{document}
